\section{Herramientas}
Para el desarrollo de la aplicación se han utilizado diversas herramientas y tecnologías que facilitan la 
implementación, gestión y despliegue del proyecto. A continuación se describen las principales herramientas 
empleadas:
\subsection{Lenguajes de Programación:}
\begin{itemize}
    \item \textbf{Python:} Utilizado para el desarrollo del backend de la aplicación, aprovechando frameworks como 
    Flask FlaskWebSockets para la creación de APIs RESTful.
    \item \textbf{Kotlin:} Empleado para el desarrollo de la aplicación móvil en Android.
    \item \textbf{Docker:} Utilizado para la containerización y despliegue de la aplicación, 
    facilitando la gestión de entornos y dependencias.
    \item \textbf{SQL:} Lenguaje utilizado para la gestión y manipulación de la base de datos PostgreSQL.
\end{itemize}
\subsection{Herramientas de Desarrollo:}
\begin{itemize}
    \item \textbf{Android Studio:} Entorno de desarrollo integrado (IDE) utilizado para el desarrollo de la 
    aplicación móvil en Android.
    \item \textbf{Visual Studio Code:} IDE utilizado para el desarrollo del backend y la edición de archivos 
    relacionados con el proyecto.
    \item \textbf{Curl:} Herramienta utilizada para probar y depurar las APIs RESTful del backend.
    \item \textbf{Bash:} Utilizado para la automatización de tareas mediante scripts.
    \item \textbf{Git y GitHub:} Sistema de control de versiones y plataforma de alojamiento de código fuente 
    utilizada para gestionar el desarrollo colaborativo del proyecto.
    \item \textbf{LaTeX:} Utilizado para la redacción del informe y la documentación del proyecto.
    \item \textbf{Figma:} Herramienta de diseño utilizada para la creación de prototipos y diseños de la interfaz.
    \item \textbf{AWS CloudFormation:} Utilizado para la definición y despliegue de la infraestructura en la nube.
    \item \textbf{Docker:} Herramienta para definir y ejecutar aplicaciones Docker de múltiples contenedores.
    \item \textbf{Trello:} Herramienta de gestión de proyectos utilizada para organizar tareas y facilitar la colaboración en el equipo.
    \item \textbf{GitHub Actions:} Utilizado para la integración continua y despliegue continuo (CI/CD) del proyecto.
\end{itemize}